\documentclass[11pt]{article}
\usepackage[margin=1in]{geometry}
\usepackage{amsmath,amssymb,amsthm}
\usepackage{graphicx}
\usepackage[hidelinks]{hyperref}

\newtheorem{theorem}{Theorem}
\newtheorem{lemma}{Lemma}
\newtheorem{definition}{Definition}
\newtheorem{remark}{Remark}

\newcommand{\R}{\mathbb R}
\newcommand{\Q}{\mathbb Q}
\newcommand{\N}{\mathbb N}
\newcommand{\eps}{\varepsilon}

\title{A Compact Ordered Metric Space Which Is Not L-Ordered}
\author{Run \texttt{fa\_banach\_001}, lane 17}
\date{June 19, 2026}

\begin{document}
\maketitle

\section*{Source Problem}

The source paper is T. Fritz and P. Perrone,
\emph{Stochastic order on metric spaces and the ordered Kantorovich monad},
arXiv:1808.09898.  In Section 4.1, after introducing L-ordered spaces, the
authors ask:

\begin{quote}
Problem 4.1.3. Is every compact ordered metric space automatically
L-ordered?
\end{quote}

The problem appears on page 18 of the arXiv PDF.

\begin{center}
\includegraphics[width=\textwidth]{figures/open_problem_crop.png}
\end{center}

\section*{Definitions}

Following the source paper, an ordered metric space is a metric space equipped
with a partial order whose graph is closed in the product metric topology.  It
is L-ordered if, for every $x,y\in X$,
\[
  x\le y
  \quad\Longleftrightarrow\quad
  f(x)\le f(y)
  \text{ for every short monotone } f:X\to\R .
\]
Here short means $1$-Lipschitz.

\section*{Proof Intuition}

Compactness prevents the source paper's noncompact counterexample from being
copied directly: if a short metric bridge collapses to a genuine limit point,
closedness of the order tends to create new order relations.  The example
below avoids this by placing many finite ``almost chains'' inside the compact
interval.  Each link is an actual order relation followed by a tiny metric
gap; the total metric error in the $n$th chain tends to zero, so every short
monotone real-valued function is forced to satisfy $f(0)\le f(1)$.

At the same time, the order relation itself contains no composable nontrivial
pairs.  All nontrivial lower endpoints are rational and all nontrivial upper
endpoints are irrational.  Thus the order is transitive for a simple discrete
reason, while closedness holds because the off-diagonal pairs approach only
the diagonal.

\begin{theorem}
There exists a compact ordered metric space which is not L-ordered.  In fact,
the usual compact metric space $[0,1]$ admits such a closed partial order.
\end{theorem}

\begin{proof}
Let $X=[0,1]$ with its usual metric.  For every integer $n\ge2$, put
\[
  \eps_n=\frac{\sqrt2}{n^3}.
\]
For $0\le i\le n-1$, define
\[
  a_{n,i}=\frac{i}{n},
  \qquad
  b_{n,i}=\frac{i+1}{n}-\eps_n .
\]
Since $0<\eps_n<1/n$, all $b_{n,i}$ lie in $[0,1]$.  Each $a_{n,i}$ is
rational, while each $b_{n,i}$ is irrational.

Define a relation $\le$ on $X$ by declaring
\[
  x\le y
  \quad\Longleftrightarrow\quad
  x=y
  \text{ or }
  (x,y)=(a_{n,i},b_{n,i})
  \text{ for some } n\ge2,\ 0\le i\le n-1 .
\]
This relation is reflexive by construction.  It is antisymmetric because a
nontrivial related pair has rational first coordinate and irrational second
coordinate, so the reverse pair cannot also be nontrivial and equality is
impossible.

It is transitive as well.  Suppose $x\le y$ and $y\le z$.  If either relation
is equality, the conclusion is immediate.  Otherwise, $x\le y$ is a
nontrivial pair, so $y$ is irrational; but every first coordinate of a
nontrivial pair is rational.  Hence $y\le z$ cannot be nontrivial, a
contradiction to the remaining case.  Therefore $\le$ is a partial order.

We next prove that its graph is closed.  Let
\[
  E=\{(a_{n,i},b_{n,i}): n\ge2,\ 0\le i\le n-1\}.
\]
The graph is $\Delta_X\cup E$, where $\Delta_X$ is the diagonal.  Consider a
convergent sequence in $E$, say
\[
  (a_{n_k,i_k},b_{n_k,i_k})\to (u,v).
\]
If the integers $n_k$ are bounded, then only finitely many points of $E$ can
occur, and a constant subsequence shows that $(u,v)\in E$.  If
$n_k\to\infty$, then
\[
  |b_{n_k,i_k}-a_{n_k,i_k}|
  =\frac1{n_k}-\eps_{n_k}\longrightarrow 0,
\]
so $u=v$ and $(u,v)\in\Delta_X$.  Thus $\Delta_X\cup E$ is closed in
$X\times X$.  Hence $X$ is a compact ordered metric space.

We now show that $X$ is not L-ordered.  Let $f:X\to\R$ be any short monotone
map.  Fix $n\ge2$.  Monotonicity gives
\[
  f(a_{n,i})\le f(b_{n,i})
\]
for all $0\le i\le n-1$.  Since $f$ is short and
$|b_{n,i}-a_{n,i+1}|=\eps_n$, we also have
\[
  f(b_{n,i})\le f(a_{n,i+1})+\eps_n .
\]
Combining these inequalities yields
\[
  f(a_{n,i})\le f(a_{n,i+1})+\eps_n.
\]
Iterating from $i=0$ to $i=n-1$ gives
\[
  f(0)=f(a_{n,0})
  \le f(a_{n,n})+n\eps_n
  = f(1)+\frac{\sqrt2}{n^2}.
\]
Letting $n\to\infty$ gives $f(0)\le f(1)$ for every short monotone
$f:X\to\R$.

However, $0\nleq 1$ in the order: the pair $(0,1)$ is neither diagonal nor
one of the pairs $(a_{n,i},b_{n,i})$, since all such second coordinates are
strictly less than $1$ when $i=n-1$ and are even smaller otherwise.  Thus
short monotone functions fail to detect the order, and $X$ is not
L-ordered.
\end{proof}

\section*{Verification Notes}

The proof is symbolic and uses no computational search.  The key verifications
are:
\begin{itemize}
\item $\eps_n=\sqrt2/n^3$ is irrational and satisfies $0<\eps_n<1/n$ for
  $n\ge2$.
\item Nontrivial first coordinates are rational, while nontrivial second
  coordinates are irrational, so no two nontrivial order pairs compose.
\item The only possible new accumulation points of nontrivial order pairs are
  diagonal points, because their coordinate gaps are $1/n-\eps_n\to0$.
\item The finite-chain estimate
  $f(0)\le f(1)+n\eps_n=f(1)+\sqrt2/n^2$ holds for every short monotone
  $f$.
\end{itemize}

No simple computational contradiction is possible here: the construction is
explicit, compactness is inherited from the usual interval, and all order
checks are contained in the preceding proof.

\section*{Novelty Check and Limitations}

The local run indexes were searched for arXiv:1808.09898, the source title,
and phrases including \texttt{compact ordered metric space},
\texttt{L-ordered}, \texttt{short monotone}, and
\texttt{ordered Kantorovich}.  No existing solution packet or proof-gap entry
for this question was found.  The local arXiv source corpus was searched for
the same phrases; it found the source paper and related terminology but no
later answer to Problem 4.1.3.

External web searches were also run for the exact sentence ``Is every compact
ordered metric space automatically L-ordered'' and close variants involving
L-ordered metric spaces, compact ordered metric spaces, counterexamples, and
short monotone maps.  These searches returned the source paper but no explicit
later solution.  This is a bounded novelty check, not a publication-level
exhaustive literature review.

The counterexample answers the exact problem as stated in the source paper.
It uses a highly nonstandard closed order on the standard compact metric
interval; it does not contradict the source paper's positive classes such as
subsets of ordered Banach spaces.

\section*{Human-Review Recommendation}

This should be reviewed as a likely valid counterexample.  The reviewer should
focus on the closedness and transitivity of the sparse relation, and on the
iteration of the Lipschitz bridges from $0$ to $1$.

\begin{thebibliography}{9}
\bibitem{FP}
T. Fritz and P. Perrone,
\emph{Stochastic order on metric spaces and the ordered Kantorovich monad},
arXiv:1808.09898.
\end{thebibliography}

\end{document}
